% ============================================================
% ch06.tex —— 第 6 章 费马小定理与第一把密码锁
% ============================================================
\chapter{费马小定理与第一把密码锁}

这一章要做两件事。第一件：发现余数王国里一条深刻得反常的定律——素数的世界里，幂会准时归一。第二件：用它在十分钟内造出人类历史上最重要的锁。造锁的部分只需要小学乘法，你随时可以停下书来跟家人现场表演。

\begin{kaiwei}
不许真的算出 $2^{100}$（一个 $31$ 位数），你能知道它除以 $11$ 余几吗？更进一步：$3^{100}$ 除以 $7$？$2^{10^{6}}$ 除以 $101$？
\end{kaiwei}

\section{一个漂亮得可疑的规律}

工具还是第 1 章的余数替换，边乘边压。先摸底，取素数 $11$，底数 $2$：
\[
2^{10} = 1024 = 93\times 11 + 1 \;\Longrightarrow\; 2^{10} \equiv 1 \pmod{11}.
\]
余数恰好是 $1$。只是碰巧吗？换一组，素数 $5$，底数 $3$：
\[
3^4 = 81 = 16\times5 + 1 \;\Longrightarrow\; 3^4 \equiv 1 \pmod 5 .
\]
再换，素数 $7$，底数 $2$：
\[
2^{6} = 64 = 9\times7 + 1 \;\Longrightarrow\; 2^{6} \equiv 1 \pmod 7 .
\]
三组全中，而且注意指数与模数的关系：$10 = 11 - 1$，$4 = 5-1$，$6 = 7-1$——\textbf{指数总是"素数减一"}。三百多年前（1640 年），费马把这封信写给朋友弗雷尼克勒，信末照例加了他那句名言式的免责声明："要不是怕太长，我早就把证明给你了。"（他一生经常这样——留谜面不留谜底，其中最著名的一个谜"页边太窄写不下"让全人类忙了 358 年，见第 32 章。）

\begin{dingli}[费马小定理]
$p$ 是素数，$a$ 是不被 $p$ 整除的整数，则
\[
a^{\,p-1} \equiv 1 \pmod{p}.
\]
\end{dingli}

\textbf{为什么偏偏是素数？}下面给一个能看到骨架的直觉版证明。舞台：$Z_p$ 去掉 $0$ 的世界（居民 $1, 2, \dots, p-1$，共 $p-1$ 个非零房间），运算取模 $p$ 乘法。取一名居民 $a$，让它的幂排队出发：
\[
a^1,\; a^2,\; a^3,\; \dots \pmod p .
\]
每乘一次 $a$ 就搬一次家。世界只有 $p-1$ 个房间，\textbf{迟早撞房}：存在 $i < j$ 使 $a^i \equiv a^j$。关键一步来了——把两边的 $a^i$ \textbf{约掉}（合法性稍后说明），得
\[
a^{j-i} \equiv 1 \pmod p .
\]
也就是说：撞房不发生在别处，\textbf{恰恰发生在"$1$"这个房间}——因为一旦中途出现 $a^{j-i} \equiv 1$，队伍就开始循环。设第一次回到 $1$ 用了 $d$ 步（$a^d \equiv 1$），则整个队伍沿着 $d$ 个房间转圈：$a, a^2, \dots, a^d = 1, a, a^2, \cdots$。

现在数人头。考虑"由 $a$ 的幂组成的集合"$H = \{1, a, a^2, \dots, a^{d-1}\}$——这 $d$ 个成员自成一国（幂乘幂还是幂，封闭；每个成员的逆也在队伍里：$a^{k}$ 的逆是 $a^{d-k}$，因为 $a^k\cdot a^{d-k} = a^d \equiv 1$）。再把任意居民 $b$ 拉来，做"平移"$bH = \{b, ba, ba^2, \dots\}$——又是 $d$ 个成员（互不相同，因为若 $ba^i = ba^j$，约掉 $b$ 又见 $a^{j-i}\equiv 1$，与"最小周期 $d$"矛盾）。这些"平移副本"要么与 $H$ 完全重合，要么完全不相交（共享一个成员就能推出全体重合——代数上倒一步就到）。于是 $p-1$ 个居民被切成若干个大小恰为 $d$ 的块：
\[
p - 1 = d \times (\text{块数}) .
\]
$d$ 整除 $p-1$！设 $p-1 = dk$，则
\[
a^{p-1} = (a^d)^k \equiv 1^k = 1 \pmod p .
\]
证毕。（"约掉 $a^i$ 或 $b$"这一步的完全清白证明需要模素数世界里逆元存在——第 10 章发证。眼下先记住：素数的世界里"约去"永远合法，这是它区别于合数世界的特权。）

你在这一段里其实已经提前看见了第 9 章的主定理（拉格朗日定理：\textbf{子群的大小整除群的大小}）的完整演出——$H$ 是子群，$bH$ 是陪集，"切块"就是切蛋糕。数学的剧本常常这样：先在一个具体剧场（模 $p$ 世界）把戏演一遍，观众看懂了，再宣布"刚才那套其实适用于所有剧场"。第 9 章我们就是这么干的。

\section{驯服指数：快速幂}

立刻兑现开胃问题。$2^{100} \bmod 11$：由 $2^{10} \equiv 1 \pmod{11}$，
\[
2^{100} = (2^{10})^{10} \equiv 1^{10} = 1 \pmod{11}.
\]
一行收工。再看 $3^{100} \bmod 7$：$3^6 \equiv 1$（费马小定理），$100 = 6\times16 + 4$，所以
\[
3^{100} = (3^6)^{16} \times 3^4 \equiv 1 \times 81 = 81 = 11\times7 + 4 \equiv 4 \pmod 7 .
\]
套路总结成三步：\textbf{（1）用费马小定理把指数的"周期"定为 $p-1$；（2）把大指数除以 $p-1$，只留余数 $r$；（3）计算 $a^r \bmod p$（这步边乘边压，数始终很小）}。指数再大——一万位也行——都逃不出手掌心，因为决定余数的只是指数除以 $p-1$ 的\textbf{余数}。

这台机器有个正式名字：\textbf{快速幂取模}。它是现代密码学的发动机零件，第 32 章椭圆曲线密码里还要再用一次。

\section{造锁：全世界最重要的十分钟}

先交代需求。你想让全世界都能给你寄\textbf{加密信}，但只有你能读。对称的加密（比如约定"每个字母后移三位"）要求你和每个朋友先秘密碰头交换口令——互联网上有几十亿陌生人，碰不过来。理想方案叫\textbf{公钥密码}：
\begin{itemize}
\item \textbf{锁}（加密方法）：公开挂出，任何人都能拿来把消息锁上；
\item \textbf{钥匙}（解密方法）：只有你持有；
\item 关键安全性要求：\textbf{看着锁，造不出钥匙}。
\end{itemize}
1977 年，Rivest、Shamir、Adleman 三人（首字母即 RSA）发现：\textbf{费马小定理 + "大数分解极难"（第 3、4 章的老朋友）恰好能造出这把锁}。用玩具尺寸走一遍全流程——你现在手里的知识一步不缺。

\paragraph{第一步：造锁（接收方做）。}挑两个素数，$p = 3$，$q = 11$。相乘：
\[
N = p\times q = 33 .
\]
再算一个藏在幕后的数：$(p-1)(q-1) = 2\times10 = 20$（它叫 $\varphi$，暂时当它是"锁的齿轮数"，出处第 9 章再讲）。挑一个与 $20$ 互素的指数：$e = 3$（$\gcd(3,20)=1$ \checkmark）。\textbf{把 $(N, e) = (33, 3)$ 公开广播}——这就是锁，挂在网络上谁都能抄走。

\paragraph{第二步：配钥匙（接收方做，保密）。}找一个 $d$ 使
\[
e\times d \equiv 1 \pmod{(p-1)(q-1)} \quad\text{即}\quad 3d \equiv 1 \pmod{20}.
\]
试：$3\times7 = 21 \equiv 1$ \checkmark，所以 $d = 7$。\textbf{钥匙是 $7$}，连同"$p=3, q=11$"的拆分方式一起锁进保险柜。

\paragraph{第三步：朋友加密。}朋友要把消息 $m = 4$ 寄给你。他用公开的锁 $(33, 3)$：
\[
c = m^e \bmod N = 4^3 \bmod 33 = 64 \bmod 33 = 31 .
\]
（$64 = 33 + 31$。）他把密文 $c = 31$ 大大方方扔到公共信道上——反正是密文，谁看都看不懂。

\paragraph{第四步：你解密。}收到 $31$，用私钥 $d = 7$：
\[
m' = c^d \bmod N = 31^7 \bmod 33 .
\]
别慌，用第 1 章的手艺\textbf{边乘边压}，并且先做个聪明替换：$31 = 33 - 2$，所以 $31 \equiv -2 \pmod{33}$。于是
\[
31^7 \equiv (-2)^7 = -128 \pmod{33}.
\]
再算 $-128$ 模 $33$：$128 = 3\times33 + 29 = 99 + 29$，故 $-128 \equiv -29 \equiv 33 - 29 = 4 \pmod{33}$。
\[
m' = 4 \;\checkmark
\]
原始消息完好归来。

\paragraph{偷窥者的困境。}他截到了 $c = 31$，也查得到公开的锁 $(33, 3)$。要还原消息，他要么解"哪个数的立方除以 $33$ 余 $31$"（玩具尺寸下试遍 $33$ 个数能破）；要么——这是 RSA 真正的命门——把 $N = 33$ \textbf{分解}回 $3\times11$，自己算出齿轮数 $20$ 和私钥 $7$。玩具锁当然一掰就开；但把 $N$ 换成两个三百位素数的乘积（六百多位的大数），分解它所需的计算量按已知最好算法也是天文数字——\textbf{用全世界的计算机联机跑，时间长过宇宙年龄}。而"造锁"只需两个三百位素数（找素数很快，第 3 章试除法加上现代化身）、一次乘法。\textbf{乘上去一秒，拆开来一万年}——不对称性就是安全性。第 3、4 章埋的"分解大数极难"在这里正式上岗。

\begin{faming}
这套流程正式名字叫 \textbf{RSA 公钥密码系统}（Rivest--Shamir--Adleman，1978 年正式发表；英国情报机构 GCHQ 的数学家 Cocks 在 1973 年已内部发明同等系统，因保密封存至 1997 年——数学发现有时会"双胞胎式"降临）。它的数学心脏是一条我们尚未完全证明的深水区事实：
\[
m^{ed} \equiv m \pmod N \qquad \text{（对所有 } m\text{）}，
\]
其中 $ed \equiv 1 \pmod{(p-1)(q-1)}$。为什么"先 $e$ 次方再 $d$ 次方"会还原？因为 $ed$ 是"$1$ 加上齿轮数的整数倍"（$3\times7 = 21 = 1 + 20$），而费马小定理保证"幂次每过 $p-1$（或 $q-1$）归零重置"。完整链条（对 $p$、$q$ 两个世界分别用费马小定理，再用中国剩余定理拼回来）在第 9、10 章补上最后一颗螺丝。
\end{faming}

\begin{lishi}
RSA 三人组的原始实验用信件是莎士比亚式的挑逗："THE MAGIC WORDS ARE SQUEAMISH OSSIFRAGE"（魔咒是"挑剔的鱼鹰"）。1994 年，全球六百名志愿者用业余计算机花了八个月，暴力分解了 RSA-129（129 位十进制的 $N$，1977 年《科学美国人》专栏悬赏 100 美元），解出这条消息——这次破译不是攻击失效，恰恰是它坚挺的证明：破 129 位要八个月，而实用密钥早已升到六百位以上。你每次看到浏览器地址栏的小锁图标，背后都有"费马小定理 + 大数分解之难"在站岗：公元前 300 年的欧几里得、公元 4 世纪的《孙子算经》、17 世纪的费马，隔空联名为你的网购付款签字。
\end{lishi}

\begin{dongshou}
\begin{itemize}
  \yisi 用同一把锁 $(N=33, e=3)$ 加密 $m = 2$，得密文 $c$；再用 $d = 7$ 解密验证。（$2^3 = 8$，$c = 8$；$8^7 \bmod 33$：$8^2 = 64 \equiv 31 \equiv -2$，$8^4 \equiv 4$，$8^7 = 8^4\cdot8^2\cdot8 \equiv 4\times(-2)\times8 = -64 \equiv -64 + 66 = 2$ \checkmark。）
  \yier 自己造一把锁：$p = 5$，$q = 11$。算 $N$、$(p-1)(q-1)$，选个与齿轮数互素的 $e$（$3$ 或 $7$ 都行），求 $d$（解 $3d\equiv1 \pmod{40}$，试出 $d = 27$；或换 $e=7$ 解 $7d\equiv1$，$d=23$）。加密 $m = 3$，解密验证。
  \yier 证明费马小定理的"全体版"：$a^p \equiv a \pmod p$ 对\textbf{所有}整数 $a$ 成立（不要求 $p \nmid a$）。（$p\nmid a$ 时由定理两边乘 $a$ 即得；$p \mid a$ 时两边都是 $0$，也成立。）这个版本写起来更省事，是"素性测试"的常用形态。
\end{itemize}
\end{dongshou}

\begin{shaonao}
$561 = 3\times11\times17$ 有个邪门性质：对\textbf{所有}与它互素的 $a$，都有 $a^{560} \equiv 1 \pmod{561}$——可它根本不是素数！这种冒牌货叫\textbf{卡迈克尔数}（以发现者命名，最小的就是 $561$）。它们是"用费马小定理检验素数"这一方案的天然克星。请验证它的一次行骗：取 $a = 2$，证 $2^{560} \equiv 1 \pmod{561}$。思路：对三个素因数分别用费马小定理收网——$2^{2} \equiv 1 \pmod 3$，而 $560 = 280\times2$，故 $2^{560}\equiv1\pmod3$；$2^{10}\equiv1\pmod{11}$，$560 = 56\times10$，故 $\equiv 1 \pmod{11}$；$2^{16}\equiv1\pmod{17}$，$560 = 35\times16$，故 $\equiv 1 \pmod{17}$。三个世界都余 $1$，而 $3, 11, 17$ 两两互素，用第 5 章的中国剩余定理拼回：$2^{560}\equiv1\pmod{561}$ \checkmark。你刚刚亲手拆穿了一场"完美犯罪"——也顺便看到了 CRT 在数论里的日常用法。
\end{shaonao}

\begin{chengshi}
直觉版证明里"幂的集合切块、每块大小为 $d$、$d$ 整除 $p-1$"即拉格朗日定理的特例，第 9 章给完整版；"$ed$ 还原"的完整证明（CRT 拼合两个世界）在第 9、10 章之间补齐。另有一条重要提醒：真实世界的 RSA 还需要\textbf{填充方案}（直接加密小数字会被"小指数攻击"等手段破解），本书演示的是裸数学内核。最后，RSA 的安全性建立在一个\textbf{从未被证明的假设}上（"分解没有快速算法"）——它也许某天被证伪，这就是密码学"建议密钥长度逐年加码"的原因。数学的地毯下面，永远还有下一层。
\end{chengshi}

篇一收工。清点行李：余数王国（1）、裁边角机器与裴蜀等式（2）、原子理论（3）、反证法与素数荒漠（4）、中国剩余定理（5）、费马小定理与密码锁（6）。下一篇，我们退后一步，眯起眼睛看这六件兵器——你会发现它们身上长着同一副骨架。看穿骨架的能力，就是"抽象"二字的全部含义。
